% =====================================================================
% ARQUIVO TEMPORÁRIO DE RASCUNHO
%
% Conteúdo: tradução e ampliação da dedução de Sheldon Ross (A First
% Course in Probability, 8a ed., Seção 5.5, pp. 210-211), na qual a
% falta de memória caracteriza a distribuição exponencial.
%
% Destino previsto em vet-continuos.tex: inserir em torno da linha 1035,
% logo após o cabeçalho da subseção
%   \subsection{Corridas exponenciais e amostragem softmax}
% e antes de \begin{example}[Relógios exponenciais, ...].
% Ao inserir no documento principal, troque \section por \subsection
% (ou ajuste a hierarquia) conforme a posição escolhida.
%
% Este arquivo é autônomo apenas para revisão e compilação isolada.
% =====================================================================
% !TeX spellcheck = pt_BR
\documentclass[12pt,portuguese]{article}

%----------------------------------------------------------------------
% PACOTES ESSENCIAIS E TRADUÇÃO
%----------------------------------------------------------------------
\usepackage[brazilian]{babel}     % Suporte à língua portuguesa
\usepackage[utf8]{inputenc}       % Codificação de entrada
\usepackage[T1]{fontenc}          % Codificação da fonte
\usepackage{lmodern}              % Fonte escalável

%----------------------------------------------------------------------
% PACOTES MATEMÁTICOS
%----------------------------------------------------------------------
\usepackage{amsmath, amsthm, amssymb, amstext, mathtools}
\usepackage{mathrsfs}             % Para fontes como \mathscr
\usepackage{dsfont}               % Para o indicador de conjunto \mathds{1}
\usepackage[all]{xy}              % Para diagramas comutativos
\usepackage{imakeidx}             % Criação de Lista de Símbolos / Índice

% Desigualdades sempre com símbolos inclinados
\let\leq\leqslant
\let\le\leqslant
\let\ge\geqslant
\let\geq\geqslant

%----------------------------------------------------------------------
% FORMATAÇÃO, LAYOUT E CORES
%----------------------------------------------------------------------
\usepackage{geometry}
\usepackage{varwidth}
\usepackage{graphicx}
\usepackage[pdftex]{xcolor}
\usepackage{titlesec}
\usepackage{appendix}
\usepackage{enumitem}
\usepackage{float}
\usepackage{booktabs}
\usepackage{array}
\usepackage{colortbl}
\usepackage{caption}
\usepackage{stackengine}
\captionsetup{width=0.8\textwidth}

% Definição de Cores
\definecolor{teal}{RGB}{15,118,110}
\definecolor{text}{RGB}{28,25,23}
\definecolor{muted}{RGB}{87,83,78}
\definecolor{card}{RGB}{255,255,255}
\definecolor{border}{RGB}{231,229,228}

% Definição de margens e espaçamento
\setlength{\topmargin}{-0.7in}
\setlength{\oddsidemargin}{0in}
\setlength{\textheight}{9.55in}
\setlength{\textwidth}{6.5in}
\setlength{\baselineskip}{12mm}

% Formatação dos títulos de seção
\titleformat{\section}{\Large\bfseries}{\thesection.}{0.5em}{}
\titleformat{\subsection}{\large\bfseries}{\thesubsection.}{0.5em}{}
\titleformat{\subsubsection}{\normalfont\itshape}{\thesubsubsection.}{0.6em}{}

%----------------------------------------------------------------------
% TIKZ E GRÁFICOS
%----------------------------------------------------------------------
\usepackage{tikz}
\usepackage{pgfplots}
\pgfplotsset{compat=1.17}
\usetikzlibrary{positioning,calc,fit,nfold,decorations.pathreplacing,calligraphy,tikzmark,patterns,arrows.meta,bending,hobby}

%----------------------------------------------------------------------
% REFERÊNCIAS, LINKS E METADADOS DO PDF
%----------------------------------------------------------------------
\usepackage[pdftex, colorlinks, linkcolor={red!50!black},
citecolor={red}, urlcolor={blue!80!black}]{hyperref}

\hypersetup{
pdftitle={A distribuição exponencial e a falta de memória},
pdfauthor={L. Cioletti},
pdfsubject={Caracterização da distribuição exponencial pela falta de memória},
pdfkeywords={distribuição exponencial, falta de memória, equação funcional, MAT-UnB}
}

%----------------------------------------------------------------------
% ALIASCNT E AMBIENTES DE TEOREMA (Cleveref)
%----------------------------------------------------------------------
\usepackage{aliascnt}
\theoremstyle{definition}
\newtheorem{theorem}{Teorema}

\newaliascnt{definition}{theorem}
\newtheorem{definition}[definition]{Definição}
\aliascntresetthe{definition}

\newaliascnt{example}{theorem}
\newtheorem{example}[example]{Exemplo}
\aliascntresetthe{example}

\newaliascnt{proposition}{theorem}
\newtheorem{proposition}[proposition]{Proposição}
\aliascntresetthe{proposition}

\newaliascnt{corollary}{theorem}
\newtheorem{corollary}[corollary]{Corolário}
\aliascntresetthe{corollary}

\newaliascnt{lemma}{theorem}
\newtheorem{lemma}[lemma]{Lema}
\aliascntresetthe{lemma}

\newaliascnt{remark}{theorem}
\newtheorem{remark}[remark]{Observação}
\aliascntresetthe{remark}

\newtheoremstyle{plainexercise}
  {3pt}{3pt}{\normalfont}{}{\bfseries}{.}{.5em}
  {\thmnumber{#2}\thmnote{ #3}}
\theoremstyle{plainexercise}
\newtheorem{exercise}{Exercício}

\theoremstyle{definition}

% Ambientes não numerados usados no texto
\newtheorem*{property}{Propriedade}
\newtheorem*{solution}{Solução}

% Se um enumerate for o primeiro conteúdo de um teorema/exercício (cabeçalho
% ainda pendente, \@inlabel true), o \leavevmode dispara o \everypar do amsthm
% que coloca o cabeçalho na própria linha, antes da lista iniciar.
\makeatletter
\setlist[enumerate]{before={\if@inlabel\leavevmode\fi}}
\makeatother

% Nomes em português para os ambientes
\AtBeginDocument{\renewcommand{\proofname}{\normalfont\textbf{Prova}}}
\renewcommand{\qedsymbol}{$\blacksquare$}

% Cleveref (deve vir depois de hyperref e amsthm)
\usepackage[brazilian,nameinlink]{cleveref}

\crefname{theorem}{teorema}{teoremas}
\Crefname{theorem}{Teorema}{Teoremas}
\crefname{definition}{definição}{definições}
\Crefname{definition}{Definição}{Definições}
\crefname{lemma}{lema}{lemas}
\Crefname{lemma}{Lema}{Lemas}
\crefname{proposition}{proposição}{proposições}
\Crefname{proposition}{Proposição}{Proposições}
\crefname{corollary}{corolário}{corolários}
\Crefname{corollary}{Corolário}{Corolários}
\crefname{example}{exemplo}{exemplos}
\Crefname{example}{Exemplo}{Exemplos}
\crefname{exercise}{exercício}{exercícios}
\Crefname{exercise}{Exercício}{Exercícios}
\crefname{remark}{observação}{observações}
\Crefname{remark}{Observação}{Observações}
\crefname{section}{seção}{seções}
\Crefname{section}{Seção}{Seções}

%----------------------------------------------------------------------
% MACROS ADICIONAIS
%----------------------------------------------------------------------
\let\oldemptyset\emptyset
\let\emptyset\varnothing
\graphicspath{{./}{../arquivos-aux/}}

% Seno na convenção em português
\renewcommand{\sin}{\operatorname{sen}}

%----------------------------------------------------------------------
% DOCUMENTO
%----------------------------------------------------------------------
\begin{document}

\begin{center}
{\LARGE\bf A distribuição exponencial e a falta de memória}
\vspace{0.3cm}

{\large\bf Rascunho temporário para revisão}
\vspace{0.3cm}

{\footnotesize Outubro de 2026}
\end{center}

\vspace{0.4cm}

% =====================================================================
% SEÇÃO A INSERIR EM vet-continuos.tex (por volta da linha 1035)
% =====================================================================

\section{A distribuição exponencial e a falta de memória}
\label{sec:exponencial-falta-memoria}

Nas aplicações, muitas variáveis aleatórias representam o \emph{tempo de
espera} até a ocorrência de um evento. O tempo até a chegada de uma chamada
telefônica, o tempo até o decaimento de um núcleo radioativo e, como veremos
na corrida de relógios da subseção seguinte, o tempo até o primeiro relógio
tocar são exemplos típicos. Em todos esses casos costumamos supor que o
processo \emph{não envelhece}: a probabilidade de o evento ocorrer em um
intervalo curto é proporcional ao comprimento do intervalo e não depende do
que aconteceu no passado.

Nesta seção vamos mostrar que essa única exigência determina a forma da
distribuição. Em outras palavras, a exponencial não é apenas uma escolha
conveniente para modelar tempos de espera: ela é a \emph{única} lei de
probabilidade compatível com a hipótese de que o processo não envelhece.
A exposição acompanha a apresentação clássica de Sheldon Ross em
\emph{A First Course in Probability}.

Organizamos o argumento em três partes. Primeiro definimos a distribuição
exponencial e calculamos sua função de distribuição. Em seguida introduzimos
a propriedade de falta de memória, mostramos que a exponencial possui essa
propriedade e provamos a recíproca. Por fim, registramos duas consequências: a
interpretação de $\lambda$ como taxa de falha constante e a obtenção da
exponencial como tempo de espera em um processo de eventos raros.

\subsection{Definição e função de distribuição}

\begin{definition}[Distribuição exponencial]\label{def:distribuicao-exponencial}
Seja $\lambda>0$. Dizemos que uma variável aleatória $X$ tem
\emph{distribuição exponencial de parâmetro} $\lambda$, e escrevemos
$X\sim\mathrm{Exp}(\lambda)$, quando admite a função densidade
\[
  f_X(x)=
  \begin{cases}
    \lambda e^{-\lambda x}, & \text{se } x\geqslant 0,
    \\[0.2cm]
    0, & \text{se } x<0.
  \end{cases}
\]
O número $\lambda$ também é chamado de \emph{taxa} da distribuição.
\end{definition}

Antes de usar essa densidade, verificamos que ela define de fato uma
probabilidade. A função é não negativa por construção, e a normalização segue
do cálculo direto
\begin{align*}
  \int_{-\infty}^{\infty} f_X(x)\,dx
  &= \int_0^{\infty} \lambda e^{-\lambda x}\,dx
  \\[0.3cm]
  &= -e^{-\lambda x}\Big|_0^{\infty}
  = 1.
\end{align*}
Portanto $f_X$ é uma densidade. Integrando de $0$ a $a\geqslant 0$, obtemos a
função de distribuição
\begin{equation}\label{eq:fda-exponencial}
  F_X(a)=\mathbb{P}(X\leqslant a)
  =\int_0^a \lambda e^{-\lambda x}\,dx
  =1-e^{-\lambda a}.
\end{equation}
É conveniente registrar a \emph{cauda} dessa distribuição, isto é, a
probabilidade de $X$ exceder um valor dado. Para $a\geqslant 0$,
\begin{equation}\label{eq:cauda-exponencial}
  G(a)\equiv\mathbb{P}(X>a)=1-F_X(a)=e^{-\lambda a}.
\end{equation}
Note que $G$ é não crescente e contínua, com $G(0)=1$.

\begin{figure}[htbp]
\centering
\begin{tikzpicture}
\begin{axis}[
  width=0.78\textwidth, height=6.2cm,
  axis lines=left,
  xlabel={$x$}, ylabel={$f_X(x)$},
  xmin=0, xmax=8, ymin=0, ymax=2.1,
  domain=0:8, samples=200,
  legend style={at={(0.97,0.97)},anchor=north east,draw=none,fill=none},
  every axis plot/.append style={line width=1.1pt},
]
\addplot[teal] {0.5*exp(-0.5*x)};
\addplot[muted] {1.0*exp(-1.0*x)};
\addplot[black] {2.0*exp(-2.0*x)};
\legend{$\lambda=0{,}5$,$\lambda=1$,$\lambda=2$}
\end{axis}
\end{tikzpicture}
\caption{Densidades exponenciais para três taxas. Aumentar $\lambda$
concentra a massa perto de $0$ e reduz o tempo médio de espera, que é
$1/\lambda$; a densidade em $0$ é exatamente $\lambda$.}
\label{fig:densidades-exponenciais}
\end{figure}

\begin{example}[Média, variância e o significado da taxa]\label{ex:media-variancia-exponencial}
Calculemos $\mathbb{E}[X]$ para interpretar o parâmetro $\lambda$.
Integrando por partes com $u=x$ e $dv=\lambda e^{-\lambda x}\,dx$, de modo que
$v=-e^{-\lambda x}$, obtemos
\begin{align*}
  \mathbb{E}[X]
  &=\int_0^{\infty} x\,\lambda e^{-\lambda x}\,dx
  \\[0.3cm]
  &=-x e^{-\lambda x}\Big|_0^{\infty}
    +\int_0^{\infty} e^{-\lambda x}\,dx
  \\[0.3cm]
  &=0+\frac{1}{\lambda}
  =\frac{1}{\lambda}.
\end{align*}
O termo de fronteira se anula porque $x e^{-\lambda x}\to 0$ quando
$x\to\infty$. Integrando por partes duas vezes, do mesmo modo, obtemos
$\mathbb{E}[X^2]=2/\lambda^2$ e, portanto,
\[
  \operatorname{Var}(X)
  =\mathbb{E}[X^2]-\bigl(\mathbb{E}[X]\bigr)^2
  =\frac{2}{\lambda^2}-\frac{1}{\lambda^2}
  =\frac{1}{\lambda^2}.
\]
Assim, o tempo médio de espera é $1/\lambda$. Se pensarmos em $X$ como o
tempo até o próximo evento de um processo que produz, em média, $\lambda$
eventos por unidade de tempo, então o tempo médio entre eventos é o inverso
da taxa, como esperado.
\end{example}

\subsection{A falta de memória}

Antes de enunciar a propriedade em tempo contínuo, vale lembrar o análogo
discreto, que ajuda a reconhecer de onde ela vem.

\begin{example}[A distribuição geométrica como análoga discreta]\label{ex:geometrica-sem-memoria}
Seja $N\sim\mathrm{Geom}(p)$, com $0<p<1$, o número de tentativas até o
primeiro sucesso, de modo que $\mathbb{P}(N=n)=p(1-p)^{n-1}$ para
$n\geqslant 1$. Para $k\geqslant 0$,
\[
  \mathbb{P}(N>k)=(1-p)^k,
\]
pois o evento $\{N>k\}$ equivale a não haver sucesso nas $k$ primeiras
tentativas. Logo, para $m,k\geqslant 0$,
\[
  \mathbb{P}(N>m+k\mid N>m)
  =\frac{\mathbb{P}(N>m+k)}{\mathbb{P}(N>m)}
  =\frac{(1-p)^{m+k}}{(1-p)^m}
  =(1-p)^k
  =\mathbb{P}(N>k).
\]
Ou seja, o número de tentativas já realizadas não altera a distribuição do
número de tentativas restantes. Dizemos que a geométrica é \emph{sem
memória}. Comparando a cauda $(1-p)^k$ com a cauda exponencial
$e^{-\lambda k}$, vemos que a exponencial é a versão em tempo contínuo da
mesma ideia: trocamos o decaimento discreto $(1-p)^k$ pelo decaimento
contínuo $e^{-\lambda k}$.
\end{example}

\begin{definition}[Variável aleatória sem memória]\label{def:variavel-sem-memoria}
Uma variável aleatória $X$ com $\mathbb{P}(X\geqslant 0)=1$ é \emph{sem
memória} se, para todo $s\geqslant 0$ e todo $t\geqslant 0$ com
$\mathbb{P}(X>t)>0$,
\begin{equation}\label{eq:falta-memoria-condicional}
  \mathbb{P}(X>s+t\mid X>t)=\mathbb{P}(X>s).
\end{equation}
A restrição a $\mathbb{P}(X>t)>0$ é necessária porque a probabilidade
condicional só está definida quando o evento condicionante tem probabilidade
positiva.
\end{definition}

A condição \eqref{eq:falta-memoria-condicional} compara o tempo restante de
espera após $t$ unidades com o tempo de espera inicial. Ela diz que, sabendo
que já se passou o tempo $t$ sem que o evento tenha ocorrido, a lei do tempo
\emph{adicional} é a mesma de quem começa a esperar agora. Passando a cauda
$G(x)=\mathbb{P}(X>x)$ para o lado esquerdo, vemos que a definição é
equivalente a uma equação funcional, como registramos a seguir.

\begin{proposition}[Reformulação como equação funcional]\label{prop:falta-memoria-equacao}
Uma variável aleatória não negativa $X$ é sem memória se, e somente se, sua
cauda $G(x)=\mathbb{P}(X>x)$ satisfaz
\begin{equation}\label{eq:equacao-funcional-cauda}
  G(s+t)=G(s)\,G(t)
  \qquad\text{para todos } s,t\geqslant 0.
\end{equation}
\end{proposition}
\begin{proof}
Suponha primeiro que $X$ é sem memória e fixe $s,t\geqslant 0$. Se
$G(t)>0$, a definição de probabilidade condicional e a inclusão
$\{X>s+t\}\subseteq\{X>t\}$ dão
\[
  G(s+t)=\mathbb{P}(X>s+t)
  =\mathbb{P}(X>s+t\mid X>t)\,\mathbb{P}(X>t)
  =G(s)\,G(t).
\]
Se $G(t)=0$, então, como $G$ é não crescente e $s+t\geqslant t$, temos
$G(s+t)\leqslant G(t)=0$, logo $G(s+t)=0$, e a igualdade $0=G(s)\cdot 0$
continua valendo. Reciprocamente, se \eqref{eq:equacao-funcional-cauda} vale
e $G(t)>0$, então
$\mathbb{P}(X>s)=G(s+t)/G(t)=\mathbb{P}(X>s+t\mid X>t)$.
\end{proof}

\begin{proposition}[A exponencial é sem memória]\label{prop:exponencial-sem-memoria}
Se $X\sim\mathrm{Exp}(\lambda)$, então $X$ é sem memória.
\end{proposition}
\begin{proof}
Para $X\sim\mathrm{Exp}(\lambda)$ a cauda é $G(x)=e^{-\lambda x}$, por
\eqref{eq:cauda-exponencial}. Então, para $s,t\geqslant 0$,
\[
  G(s+t)=e^{-\lambda(s+t)}=e^{-\lambda s}e^{-\lambda t}=G(s)G(t),
\]
e a afirmação segue da \Cref{prop:falta-memoria-equacao}.
\end{proof}

\subsection{A falta de memória caracteriza a exponencial}

A proposição anterior diz que toda exponencial é sem memória. Vamos agora ao
objetivo desta seção: mostrar a recíproca. Esse é o passo que explica
por que a exponencial aparece sempre que modelamos tempo de espera sem
envelhecimento.

O problema se reduz a \emph{resolver} a equação funcional
\eqref{eq:equacao-funcional-cauda}. Reconhecer a estrutura é o primeiro
passo: as funções da forma $x\mapsto e^{-\lambda x}$ multiplicam-se quando os
argumentos se somam, e são por isso os candidatos naturais. A tarefa é
mostrar que a monotonicidade e a continuidade da cauda eliminam todas as
outras soluções. Usaremos a estratégia usual para esse tipo de equação:
determinar a solução primeiro nos racionais, onde a equação é puramente
algébrica, e depois estender aos reais usando a regularidade de $G$.

\begin{theorem}[Caracterização pela falta de memória]\label{thm:caracterizacao-falta-memoria}
Seja $X$ uma variável aleatória com $\mathbb{P}(X>0)=1$. Se $X$ é sem
memória, então existe $\lambda>0$ tal que $X\sim\mathrm{Exp}(\lambda)$.
\end{theorem}
\begin{proof}
Escreva $G(x)=\mathbb{P}(X>x)$. Por hipótese,
$G(0)=\mathbb{P}(X>0)=1$. Pela \Cref{prop:falta-memoria-equacao}, vale
\begin{equation}\label{eq:eq-funcional-prova}
  G(s+t)=G(s)G(t)\qquad\text{para todos } s,t\geqslant 0.
\end{equation}

\smallskip
\noindent\textbf{Passo 1 (determinar $G$ nos racionais).}
Seja $p=G(1)$. Como $G$ é não crescente e $G(0)=1$, temos
$0\leqslant p\leqslant 1$. Vamos mostrar primeiro que $0<p<1$.

Se $p=1$, então $G(1)=1$ e, iterando \eqref{eq:eq-funcional-prova},
$G(n)=1$ para todo inteiro $n\geqslant 1$. A monotonicidade dá então
$G(x)=1$ para todo $x\geqslant 0$, o que é impossível para a cauda de uma
distribuição, pois $\mathbb{P}(X>x)\to 0$ quando $x\to\infty$. Logo $p<1$.

Se $p=0$, então $G(1)=0$. Tomando $s=t=1/2$ em
\eqref{eq:eq-funcional-prova}, obtemos $G(1/2)^2=G(1)=0$, donde
$G(1/2)=0$. Repetindo, $G(1/2^n)=0$ para todo $n$. Para qualquer $x>0$,
escolha $n$ com $1/2^n<x$; a monotonicidade dá
$G(x)\leqslant G(1/2^n)=0$, logo $G(x)=0$. Assim $\mathbb{P}(X>x)=0$ para
todo $x>0$. Como $\{X>0\}=\bigcup_{k\geqslant 1}\{X>1/k\}$ e esses eventos
crescem, a continuidade da probabilidade em sequências crescentes dá
\[
  \mathbb{P}(X>0)=\lim_{k\to\infty}\mathbb{P}(X>1/k)=0,
\]
contrariando a hipótese. Logo $p>0$.

Definimos então
\[
  \lambda=-\log p>0,
  \qquad\text{de modo que}\qquad p=e^{-\lambda}.
\]
Para inteiros $n\geqslant 1$, a iteração de \eqref{eq:eq-funcional-prova}
dá
\[
  G(n)=G(1)^n=p^n=e^{-\lambda n}.
\]
Por outro lado, aplicando \eqref{eq:eq-funcional-prova} com $s=t=1/n$,
\[
  p=G(1)=G(1/n)^n,
  \qquad\text{logo}\qquad
  G(1/n)=p^{1/n}=e^{-\lambda/n}.
\]
Combinando as duas identidades, para inteiros $m,n\geqslant 1$,
\[
  G(m/n)=G(1/n)^m=p^{m/n}=e^{-\lambda m/n}.
\]
Assim, $G(x)=e^{-\lambda x}$ para todo racional $x\geqslant 0$.

\smallskip
\noindent\textbf{Passo 2 (estender aos reais).}
Falta remover a restrição aos racionais. Usamos a continuidade à direita da
cauda. Como $G(x)=1-F_X(x)$ e toda função de distribuição é contínua à
direita, $G$ é contínua à direita. Fixe $x\geqslant 0$ e tome racionais
$x_k\downarrow x$ com $x_k\geqslant x$. Como $G(x_k)=e^{-\lambda x_k}$ e
$G$ é contínua à direita,
\[
  G(x)=\lim_{k\to\infty}G(x_k)
  =\lim_{k\to\infty}e^{-\lambda x_k}
  =e^{-\lambda x}.
\]
Portanto $\mathbb{P}(X>x)=e^{-\lambda x}$ para todo $x\geqslant 0$. Como
$\mathbb{P}(X\geqslant 0)=1$, temos $F_X(x)=0$ para $x<0$; para
$x\geqslant 0$,
\[
  F_X(x)=1-e^{-\lambda x}.
\]
Pelo Teorema Fundamental do Cálculo, $F_X$ é diferenciável em
$(0,\infty)$, com derivada $\lambda e^{-\lambda x}$. Definindo $f_X(x)=0$
para $x<0$, obtemos a densidade $f_X(x)=\lambda e^{-\lambda x}$ em todo
$\mathbb{R}$, e concluímos que $X\sim\mathrm{Exp}(\lambda)$.
\end{proof}

\begin{remark}[O caso degenerado]\label{rem:caso-degenerado}
A hipótese $\mathbb{P}(X>0)=1$ no \Cref{thm:caracterizacao-falta-memoria}
exclui o caso em que $X=0$ quase certamente. Nesse caso degenerado a
formulação condicional \eqref{eq:falta-memoria-condicional} não se aplica,
pois $\mathbb{P}(X>t)=0$ para todo $t\geqslant 0$; já a formulação
equivalente \eqref{eq:equacao-funcional-cauda} continua válida, porque seus
dois lados são nulos. Excluímos esse caso por ele não corresponder a um
tempo de espera com incerteza; a caracterização interessa-se pelo caso em
que o tempo é uma variável aleatória não trivial.
\end{remark}

\begin{remark}[Por que a regularidade é necessária]\label{rem:regularidade-equacao-funcional}
A equação $G(s+t)=G(s)G(t)$ tem, por si só, soluções patológicas: sem
qualquer hipótese de regularidade, ela admite soluções construídas a partir
de funções aditivas arbitrárias, muitas das quais não são mensuráveis. O que
separa a exponencial dessas soluções é exatamente a estrutura herdada de ser
a cauda de uma distribuição, isto é, ser não crescente e contínua à direita.
Por isso o argumento usa dois ingredientes distintos: a álgebra da equação,
no Passo 1, para obter o valor nos racionais, e a regularidade, no Passo 2,
para estender a solução. Notamos ainda que não precisamos supor que $X$ tenha
densidade nem que $G$ seja diferenciável: a existência da densidade
$f_X(x)=\lambda e^{-\lambda x}$ é uma consequência, não uma hipótese.
\end{remark}

\subsection{Duas consequências}

\subsubsection{Taxa de falha constante}

Para uma variável aleatória $X$ com densidade $f_X$ e cauda $G$, definimos a
\emph{taxa de falha} (ou taxa de risco) por
\[
  h(t)=\frac{f_X(t)}{G(t)},
  \qquad t\geqslant 0 \text{ com } G(t)>0.
\]
Para valores pequenos de $dt$, a quantidade $h(t)\,dt$ é aproximadamente a
probabilidade de $X$ cair em $(t,t+dt]$, dado que $X>t$. Dizemos que a taxa
de falha é \emph{constante} quando $h(t)$ não depende de $t$. A falta de
memória e a taxa de falha constante descrevem a mesma hipótese de
envelhecimento nulo, como mostra o cálculo a seguir.

\begin{example}[Taxa de falha da exponencial]\label{ex:taxa-falha-exponencial}
Se $X\sim\mathrm{Exp}(\lambda)$, então $f_X(t)=\lambda e^{-\lambda t}$ e
$G(t)=e^{-\lambda t}$ para $t\geqslant 0$. Logo
\[
  h(t)=\frac{\lambda e^{-\lambda t}}{e^{-\lambda t}}=\lambda,
  \qquad t\geqslant 0,
\]
isto é, a taxa de falha é constante e igual à taxa $\lambda$ da
distribuição. Reciprocamente, suponha que $X$ seja não negativa quase
certamente, que $G$ seja diferenciável e que a taxa de falha seja constante,
$h\equiv\lambda$. Então $G'(t)=-f_X(t)=-\lambda G(t)$. A única solução de
$G'=-\lambda G$ com $G(0)=1$ é $G(t)=e^{-\lambda t}$, de modo que
$X\sim\mathrm{Exp}(\lambda)$. Assim, sob a hipótese de que a taxa de falha
está definida, dizer que um objeto não envelhece, dizer que seu tempo de vida
é sem memória e dizer que sua taxa de falha é constante são três formas de
enunciar a mesma condição.
\end{example}

\subsubsection{A exponencial como tempo de espera de eventos raros}

Os tempos de espera também surgem de um segundo ponto de partida. Suponha
que eventos ocorram em instantes aleatórios, de modo que a probabilidade de
observar exatamente um evento em um intervalo muito curto seja proporcional
ao seu comprimento, que a chance de dois ou mais eventos nesse intervalo seja
desprezível, e que os eventos em intervalos disjuntos não se influenciem.
A exponencial é, então, a lei do tempo até o primeiro evento.

\begin{proposition}[Tempo até o primeiro evento]\label{prop:tempo-primeiro-evento}
Suponha que eventos ocorram de modo que, para intervalos de comprimento $h$
pequeno, a probabilidade de ocorrer exatamente um evento seja
$\lambda h+o(h)$, a de ocorrerem dois ou mais seja $o(h)$, e eventos em
intervalos disjuntos sejam independentes. Seja $T$ o tempo até o primeiro
evento. Então $T\sim\mathrm{Exp}(\lambda)$.
\end{proposition}
\begin{proof}
Fixe $t>0$ e divida $[0,t]$ em $n$ subintervalos iguais de comprimento
$t/n$. Por hipótese, a probabilidade de não ocorrer evento em um
subintervalo é
\[
  1-\bigl(\lambda t/n+o(t/n)\bigr)-o(t/n)=1-\lambda t/n-o(t/n).
\]
Pela independência, a probabilidade de não ocorrer evento em $[0,t]$ é o
produto dessas probabilidades, de modo que
\[
  \mathbb{P}(T>t)
  =\left(1-\frac{\lambda t}{n}-o\!\left(\frac{t}{n}\right)\right)^{n}
  \longrightarrow e^{-\lambda t},
\]
quando $n\to\infty$, pois $n\,o(t/n)=t\,[o(t/n)/(t/n)]\to 0$. Usamos aqui
que os eventos de não ocorrência nos subintervalos são conjuntamente
independentes e que a taxa $\lambda$ é a mesma em todos eles (homogeneidade
temporal). Esse limite é o mesmo cálculo que produz a aproximação de Poisson
pela binomial. Logo $\mathbb{P}(T>t)=e^{-\lambda t}$ e
$T\sim\mathrm{Exp}(\lambda)$.
\end{proof}

\begin{remark}[Ligação com a corrida de relógios]\label{rem:ligacao-corrida}
As duas descrições se encontram. No modelo de eventos raros, $T$ é o tempo
até o próximo evento e sua lei é exponencial. Na subseção seguinte, tomamos
diretamente $X_i\sim\mathrm{Exp}(\lambda_i)$ como o tempo até o relógio $i$
tocar. A \Cref{thm:caracterizacao-falta-memoria} mostra que essa é a única
escolha compatível com a hipótese de que cada relógio não envelhece, e a
\Cref{prop:falta-memoria-equacao} é a equação funcional que permite fatorar
as caudas no cálculo de $\mathbb{P}(T>t,J=i)$.
\end{remark}

\end{document}
